\question
Для автоматизации систем умного дома ИТМО инженерам нужно спроектировать логический модуль всего на одном типе микросхем. На складе нашлись два вида чипов, реализующих функции трех аргументов $f(x,y,z)$ и $g(x,y,z)$.

Векторы значений функций:
\[ f = (1,1,0,1,0,0,1,0) \]
\[ g = (1,0,1,1,0,0,1,0) \]

\begin{enumerate}
    \item Используя критерий Поста, проверьте, можно ли только на чипах первого типа ($f$) построить схему любой сложности (является ли полной системой).
    \item Проверьте на полноту систему чипов второго типа ($g$).
    \item Если какая-то из функций оказалась функционально полной, выразите с её помощью инверсию ($\overline{x}$) и конъюнкцию ($x \land y$), используя только суперпозицию этой функции и тождественных переменных.
\end{enumerate}

---------------

Автор -- Зеленская Дарья